\hspace*{\fill}

\section*{Originally published by}
World Scientific Publishing Co. Pte. Ltd.\\
5 Toh Tuck Link, Singapore 596224\\
USA office: 27 Warren Street, Suite 401-402, Hackensack, NJ 07601\\
UK office: 57 Shelton Street, Covent Garden, London WC2H 9HE


\section*{Open Access}
This project is open access and is licensed under the terms of the Creative Commons Attribution 4.0 International License (\href{http://creativecommons.org/licenses/by/4.0/}{http://creativecommons.org/licenses/by/4.0/}), which permits use, sharing, adaptation, distribution and reproduction in any medium or format, as long as you give appropriate credit to the original author(s) and the source, provide a link to the Creative Commons license and indicate if changes were made.

It is based on the original work published by World Scientific Publishing Co. Pte. Ltd., Singapore, 1988, ISBN 978-9971-5-0107-5 (hardcover), ISBN 978-9971-5-0996-5 (paperback), ISBN 978-981-4415-49-1 (ebook for institutions), ISBN 978-981-4578-28-8 (ebook for individuals).

\section*{Open Access funded by SCOAP (Sponsoring Consortium for Open Access Publishing in Particle Physics)}
The original ebook was converted to Open Access in 2021 through the sponsorship of SCOAP , licensed under the terms of the Creative Commons Attribution 4.0 International License (\href{http://creativecommons.org/licenses/by/4.0/}{http://creativecommons.org/licenses/by/4.0/}), which permits use, sharing, adaptation, distribution and reproduction in any medium or format, as long as you give appropriate credit to the original author(s) and the source, provide a link to the Creative Commons license and indicate if changes were made.
Original copyright © 1988 by World Scientific Publishing Co. Pte. Ltd.

\newpage

\section*{Preface}
In the years since its publication, this book has become the standard reference on the quantum theory of angular momentum.  The book provides a comprehensive treatment of the subject, covering both the theoretical foundations and practical applications of angular momentum in quantum mechanics. As an open access publication, it is freely available to researchers and students worldwide, ensuring that the knowledge contained within its pages can be accessed by anyone interested in the field. The book has been widely cited and used in research and teaching, and it continues to be a valuable resource for those studying quantum mechanics and related areas of physics.

However, there are some drawbacks to the form in which it was provided. The scanned version of the book is not searchable, and the quality of the scanned images is not always optimal. This can make it difficult to find specific information or to read the text clearly. Additionally, the scanned version may not be compatible with all devices or software, which can limit its accessibility. Also, there is no verification of the accuracy of the scanned text, which can lead to errors or misinterpretations of the content. Finally, the scanned version may not be easily editable or adaptable for different purposes, such as creating new editions or translations.

Thus this attempt to provide a more accessible and user-friendly version of the book, with improved searchability, image quality, and compatibility with different devices and software. The text has been carefully reviewed and verified for accuracy, and it is provided in a format that can be easily edited or adapted for different purposes. We hope that this version of the book will be a valuable resource for researchers, students, and anyone interested in the quantum theory of angular momentum. The author would like to thank the original publisher, World Scientific Publishing Co. Pte. Ltd., for their support in making this open access version available to the public.

Clearly in  the era of AI, I have made substantial use of AI tools to assist in the preparation of this version of the book. I have used AI tools to help with tasks such as text recognition, formatting, and editing. These tools have helped to improve the accuracy and quality of the text, as well as to streamline the process generating code to verify equations. However, I have also carefully reviewed and verified the output of these tools to ensure that the final version of the book is as accurate and reliable. Overall, I believe that the use of AI tools has been a valuable aid in the preparation of this new open access version of the book.

Of course, the author would like to acknowledge that the original work was done by D. A. Varshalovich, A. N. Moskalev, and V. K. Khersonskii, and that this new version is based on their work. The author has made every effort to ensure that the content of this new version is faithful to the original work, while also improving its accessibility and usability for modern readers.

I do no see this as a static resource--there are undoubtedly errors in the text, and I would be grateful for any corrections or suggestions for improvement. Please post these at the GitHub repository for this project at \href{}. The intention is to improve the accessibility and usability of this important work, and I welcome any feedback, and contributions, that can help achieve that goal.

Manchester, Uk, September 2026 

\section*{Original Preface}
This book deals with one of the basic topics of quantum mechanics: the theory of angular momentum and irreducible tensors. Being rather versatile, the mathematical apparatus of this theory is widely used in atomic and molecular physics, in nuclear physics and elementary particle theory. It enables one to calculate atomic, molecular and nuclear structures, energies of ground and excited states, fine and hyperfine splittings, etc. The apparatus is also very handy for evaluating the probabilities of radiative transitions, cross sections of various processes such as elastic and nonelastic scattering, different decays and reactions (both chemical and nuclear) and for studying angular distributions and polarizations of particles.

Today this apparatus is finding ever increasing use in solving practical problems relating to quantum chemistry, kinetics, plasma physics, quantum optics, radiophysics and astrophysics.

The basic ideas of the theory of angular momentum were first put forward by M. Born, P. Dirac, W. Heisenberg and W. Pauli. However, the modern version of its mathematical apparatus was developed mainly in the works of E. Wigner, J. Racah, L. Biedenharn and others who applied group theoretical methods to problems in quantum mechanics.

At present a number of good books on the theory of angular momentum have been already published. The general principles and results of the theory may be found in the books by M. Rose [31], A. Edmonds [16], U. Fano and G. Racah [18], A. P. Yutsis, I. B. Levinson and V. V. Vanagas [44], A. P. Yutsis and A. A. Bandzaitis [45], D. Brink and G. Satcher [9]. Nevertheless, many formulas and relationships essential for practical calculations have escaped these books and are either scattered in various editions, or included as appendices in papers discussing somewhat disparate topics, making them generally inaccessible. Even greater difficulties arise when one tries to use the results, as each author employs his own phase conventions, initial definitions and symbols.

The authors of this book aimed at collecting and compiling ample material on the quantum theory of angular momentum within the framework of a single system of phases and definitions. This is why, in addition to the basic theoretical results, the book also includes a great number of formulas and relationships essential for practical applications.

This edition is the translated version of our book published in the USSR in 1975. In the course of its preparation we have tried to comply with a number of suggestions from our readers. For instance, each chapter opens with a comprehensive listing of its contents to ease the search for information needed. We also included some new results relating to different aspects of angular momentum theory which have recently appeared in journals. Unfortunately the limited volume of the present book prevented us from covering all the aforementioned results. We offer sincere apologies to the authors whose results we failed to include.

The monograph is a kind of handbook. Consequently the material is presented in concise form. Most of the formulas and relationships are given without proof. Their full derivation may be found in the literature\\
listed at the end of the text. Some results which have become generally known are given without references, and for this we also apologize.

The sequence adopted is as follows: chapter, section and subsection. Many chapters are self-contained and can be read independently of the others. Sections have double numbering: the first figure denotes the number of the chapter, the second, the number of the section. Equations are numbered within the confines of the section they are included in. When referring to an equation from the same section only the number of the equation is given, e.g., (3), (27); when reference is made to an equation from another section the numbers of the chapter, section and equation are given, e.g., Eq. 4.2.(17). A similar system is adopted when referring to individual subsections, e.g., Sec. 1.2.5. For convenience the book also contains a glossary of all symbols used in the text with references to the pages where their corresponding definitions are given. The list of references is divided into parts: the first part lists books and reviews; the second, papers on different subjects; the third, tables; the fourth, references added during translation.

The authors hope that many specialists will find in the book some fresh and interesting information. The material is prepared and arranged so as to make it useful to those less familiar with theory and for students of physics. These readers can effectively use the monograph as a supplementary text to their main courses.

For those who wish to thoroughly familiarize themselves with the fundamentals of angular momentum theory we recommend the excellent new book by L. Biedenharn and J. Louck [132] Angular Momentum in Quantum Physics. Theory and Applications.

The authors wish to express their deep appreciation to D. G. Yakovlev who took the trouble of reading the English translation of the book and gave some valuable suggestions on its preparation.

Leningrad